\documentclass{article}
\usepackage[preprint,main]{neurips_2025}
\makeatletter
\renewcommand{\@noticestring}{Unsubmitted working draft; NeurIPS 2025 style used for layout only.}
\makeatother
\usepackage[utf8]{inputenc}
\usepackage[T1]{fontenc}
\usepackage{amsmath,amssymb,booktabs,microtype}
\usepackage{graphicx}
\usepackage{hyperref}
\usepackage{url}
\title{Rethinking Relational Evidence in Fixed-Candidate Retrieval:\\
Ablation Gains, Equal-Fact Alternatives, and Decision Advantage}
\author{Yapu Zhang \\
Beijing University of Technology \\
\href{mailto:zhangyapu@bjut.edu.cn}{zhangyapu@bjut.edu.cn}
\And
Xianliang Yang\thanks{Corresponding author.} \\
Microsoft \\
\href{mailto:xianya@microsoft.com}{xianya@microsoft.com}}

\begin{document}
\maketitle
\begin{abstract}
A positive graph-feature ablation does not by itself establish that graph
representation is irreplaceable or that a learned decision rule beats a simple
graph-aware rule. We separate these claims for fixed-pool STaRK-Amazon retrieval.
An 18-feature ranker on 384 held-out TRAIN queries (H) hits 278 answers, versus
264 for its fact-deprived text ablation and 263 for a candidate-shuffled graph
control, but 279 for a top-linked-entity rule using the shared inputs. A distinct
20-feature adapted ranker on 768 other TRAIN queries (I) hits 568 versus 549 and
552 for its matched controls, but 573 for the rule. H and I evaluate different
methods, not two confirmations of one method. Post hoc equal-fact text programs
reproduce every rule decision and all six original graph columns; an H-informed
structured-text learner ties the graph ranker at 278/384 on exposed H queries.
These observations support a bounded relational-feature increment over text
without relational facts, not graph-only informational value or learned-rule
superiority. A new cohort is unscored; matched end-to-end resource comparison
and human semantic review remain outstanding. This is a working draft, not an
independently confirmed submission.
\end{abstract}

\section{Introduction}
Queries about a product and a named reference may require both attribute
matching and relational evidence. A co-view or co-purchase edge can narrow a
search, but does not establish that a candidate has the requested role or is
compatible with the reference. In hybrid retrieval, an improvement relative to
an ablation that lacks those edges can be mistaken for evidence that graphs are
uniquely necessary, or that a trained graph ranker is the best decision rule.
Those claims require different controls.

We ask three questions on the same fixed 20-candidate pools: (1) does matched
relational evidence add to a ranker whose text branch has no relation facts?
(2) can a text-side procedure with the \emph{same facts} reproduce graph-driven
decisions or features? (3) does learned fusion outperform a strong, simple
rule with the same underlying candidates, rankings, reference links, and edges?
The distinction between fact availability, representation and decision policy
is the contribution of this empirical audit, not a new retrieval algorithm.
It complements studies of semi-structured retrieval \citep{stark,hybgrag,autofocus,grasp}
and graph reranking \citep{grag,subgraphrag}; neither fixed-pool reranking nor
matched ablations are claimed as firsts here.

The two frozen Amazon evaluations yield a positive answer to the first question
for their respective methods, but not to the third. Post hoc text replays
challenge a claim of fact exclusivity without constituting a matched-budget LLM
baseline. A separate exploratory MAG author task shows that an author-field
text join can reproduce graph decisions on linkable queries; it is not a
cross-domain confirmation. These distinctions matter most when a negative
strong-baseline result would otherwise be hidden behind a successful ablation.

\section{Related work and contribution boundary}
\paragraph{Retrieval from text and relations.}
STaRK \citep{stark} supplies text-and-relation entity retrieval benchmarks;
its synthetic and human query sources should not be interchanged. Our
reported Amazon H/I queries come from the official TRAIN pool, and our
fixed-pool Hit@1 cannot be compared numerically to the benchmark's
full-corpus TEST systems. HybGRAG \citep{hybgrag} and Autofocus Retrieval
\citep{autofocus} combine evidence discovery, constraints and reranking;
GRASP \citep{grasp} studies plan-guided graph retrieval and fusion. We
freeze the initial 20 products and ask which \emph{inference} follows from
an additional set of relational features at the decision stage. Such a
restriction removes important system behavior, notably candidate recall,
index updates and open-ended reasoning. It is useful for attributing one
decision change, not a replacement for those end-to-end systems.

\paragraph{Graph reranking and strong controls.}
G-RAG \citep{grag} already studies graph-based reranking of initial
retrieval results. SubgraphRAG \citep{subgraphrag} studies structural
scores and matched-candidate or random-retrieval comparisons. Matching
candidate pools, shuffling relations and applying LambdaRank
\citep{lightgbm} are therefore not method inventions. Here the
additional diagnostic is an ordering of \emph{claims}: first compare
matched graph features to a fact-deprived text arm, then ask whether
the same relation facts yield reproducible text-side actions, then
compare learned selection to an explicit graph-aware rule. Only the
first stage was an original H/I frozen ablation family; the equal-fact
parser and structured-text learner were built post hoc. Treating the
later controls as if they had been prospectively designed would erase
the very research history that motivates this distinction.

\paragraph{Statistical evidence is not deployment evidence.}
The paired-query discordance test asks whether a fixed arm yields more
benchmark fixes than harms under its null assumptions. It does not
measure semantic compatibility, marginal dollars, or robustness to a
new source of queries. The top-link rule consumes graph/link facts, so
its higher hit count challenges the particular learned fusion decision,
not the utility of graph access. Conversely, a text parser that can
recreate support moments from recorded typed edges shows that those
moments are not \emph{informationally exclusive} to a graph data
structure. It says nothing by itself about online text-model quality
or compute parity. The contribution is therefore a scoped empirical
audit of what the existing controls license one to say, not a proposed
universal replacement for relational search.

\section{Problem and evaluation roles}
Given query $q$, frozen candidate pool $C_q$ of size 20, and labeled answer
set $Y_q$, a policy outputs $\pi(q)\in C_q$. Our outcome is
\begin{equation}
  \operatorname{Hit@1}(\pi)=\frac{1}{N}\sum_{q=1}^{N}
  \mathbf{1}\{\pi(q)\in Y_q\}.
\end{equation}
The statistical unit is the query, not a candidate, graph world, LLM response,
or pair of methods. Synthetic benchmark answer sets do not certify semantic
negatives. We keep official TEST and human queries out of this study.

The graph ranker fits on 768 official TRAIN queries (A/B/F); 384 repeatedly
used TRAIN queries (C/D/E) serve development. H comprises 384 different TRAIN
queries sampled by a frozen sorted-ID permutation after method selection.
The separate adapted 20-feature method additionally fits its semantic encoder
on 1,024 disjoint TRAIN queries (S), then evaluates on 768 further TRAIN
queries (I), with its comparison family fixed before query preparation. H and I
are disjoint but differ in method and supervision; both are now exposed. Neither
is official TEST, an independent replication of the other, or an independent
cross-dataset evaluation. Query disjointness does not imply disjoint products,
graph edges, pretraining corpora, or freedom from historical research selection.

\section{Matched inputs and competing decisions}
BM25 and a frozen E5 retriever supply the RRF top-20 pool. Two LLM services
rank those same candidates; one shared reference-extraction call identifies
at most three named products per query. Each reference has at most three
frozen linking options. For interpretation $h$ and candidate $v$, let
\begin{equation}
 U_q(v,h)=\frac{1}{m_q}\sum_{j=1}^{m_q}
 \mathbf{1}\{\text{$v$ is connected to reference $j$ in $h$}\},
\end{equation}
when $m_q>0$; without a reference, all support statistics are zero. This
proxy does not enforce direction, edge-type semantics, or product compatibility.
An 18-dimensional LambdaRank fusion model \citep{lightgbm} combines six symmetric
two-rank statistics, six frozen semantic/text statistics (including a Qwen
reranker), and six uniform or link-weighted world-support statistics. Its
12-dimensional text ablation retains ranks and semantic inputs but has no
relation facts. A separate shuffled control permutes the six graph columns
across candidates within each query, using one prespecified seed for training
and evaluation. These arms share the candidate pools, available rankings,
supervision and ranker setup; the shuffle is not a randomization of the full
knowledge graph and has not been repeated across seeds.

The strongest simple comparator first chooses the highest-weight link for
each mention, then ranks candidates by its support, two-LLM RRF score and
frozen candidate order in lexicographic sequence. With no references it
reduces to RRF and the order tie-breaker. This \emph{top-link} rule has access
to the same underlying evidence, but not identical feature columns or the
same minimum deployment cost as the learned model. In particular, the learned
model does not have a hard support-first decision constraint. RRF and Borda
require fewer components than the graph and Qwen paths when deployed alone.
The adapted I method adds semantic adaptation and 20-dimensional controls;
any difference between old and adapted methods also changes training data.

\paragraph{Feature construction and ambiguity.}
Let $\mathcal M_q$ be the nonempty set of strictly valid rankings, with
$r_m(v)\in\{1,\ldots,20\}$. The six rank columns are normalized minimum,
maximum, mean and range of ranks, mean reciprocal rank and fraction of
first-place votes. These statistics are invariant to exchanging the two
services; invariance of a feature builder is not invariance of service
quality. The six nonrelational columns encode initial retrieval rank,
query/title token overlap, query/product-text token overlap, clipped text
length, a frozen Qwen yes-minus-no relevance logit, and its difference
from the query's maximum logit. No logit is treated as a calibrated
probability.

For mention $j$, define a link option set $R_{qj}$ with weights $w_{qjt}$.
Title-exact ties receive uniform weights; otherwise weights are a softmax
of frozen E5 similarities at temperature $0.05$. For
$h\in\mathcal H_q=\prod_j R_{qj}$, put
$p_q(h)=\prod_jw_{qj,h_j}$. From $U_q(v,h)$ the six graph columns are
uniform mean, minimum, maximum and standard deviation across worlds,
followed by $p_q$-weighted mean and standard deviation. At most 27 worlds
are explicitly enumerated; the product weights are a heuristic and not
calibrated entity-linking posteriors. With no mention all six columns are
zero. The support proxy treats either direction of also-buy and also-view
as a link, discarding type-specific semantics. Thus an informative feature
can still reward a wrong product role.

Each query contributes one group of 20 binary-labeled candidates to
LightGBM LambdaRank: 100 trees, learning rate $0.05$, at most seven leaves,
minimum 20 examples per leaf, $L_2$ regularization one and ranking truncation
at four. The frozen training configuration and label gain remain the
authoritative specification. Scores choose the highest candidate within
the full pool, breaking ties by frozen candidate order; the learned model
does \emph{not} restrict its decision to graph-supported items. For I, a
BGE semantic encoder is adapted on S using up to four positive and
additional negative training products per query, then frozen before
the 768-query fusion fitting. Two extra semantic columns bring the adapted
fusion to 20 features and its adapted text control to 14. Positive injection
belongs only to encoder training, never evaluation retrieval. The adapted
encoder adds supervision; old-versus-adapted scores do not isolate a single
algorithmic change.

\section{What can an ablation establish?}
\paragraph{Coverage and selection are distinct.}
For a fixed query let $O_q(D)=\mathbf{1}\{D\cap Y_q\ne\varnothing\}$
and $H_q(\pi)=\mathbf{1}\{\pi(q)\in Y_q\}$. A ranker restricted to
$C_q$ satisfies $H_q(\pi)\le O_q(C_q)$. If a graph-constrained policy
selects from a support domain $E_q\subseteq C_q$, its error decomposes as
\begin{equation}\label{eq:decomposition}
  1-H_q(\pi)=[1-O_q(C_q)]+[O_q(C_q)-O_q(E_q)]
                  +[O_q(E_q)-H_q(\pi)].
\end{equation}
The three nonnegative terms count missing candidate answers, answers excluded
by the domain, and within-domain selection errors. This is an identity for
the recorded labels, not a graph-correctness theorem. Our 18D fusion scores
\emph{all} 20 items and is not itself constrained to $E_q$; the decomposition
applies only to policies that select within their stated domain.

\paragraph{Equal facts do not imply equal computation.}
Let $F_q$ denote the recorded candidate identities, ranks, link options,
weights and edges. Both graph-feature extraction and a lossless typed textual
serialization are deterministic functions of $F_q$. If a parser preserves
those fields, composing it with the support computation reproduces the same
support values and any deterministic rule using them. This is a constructive
\emph{representational} observation, not an assertion that a bounded-token
language model will learn the parser, that its inputs fit a fixed context,
or that text indexing and graph traversal have equal latency or cost.
Fitting a separate lossy text learner is therefore a distinct empirical test.

\paragraph{Decision opportunity is label blind.}
For two fixed policies $g,t$ on $Q$, put
$D(g,t)=\{q\in Q:g(q)\ne t(q)\}$ and let $F$ and $B$ count queries
where $g$ alone and $t$ alone hit, respectively. Directly from the
definitions,
\begin{equation}\label{eq:disagreement}
 \operatorname{Hit@1}(g)-\operatorname{Hit@1}(t)
 =\frac{F-B}{|Q|},\qquad
 \left|\operatorname{Hit@1}(g)-\operatorname{Hit@1}(t)\right|
 \leq \frac{|D(g,t)|}{|Q|}.
\end{equation}
Agreement in action forces agreement in hit status for \emph{every} possible
answer set; disagreement does not guarantee a different hit status. A
label-blind minimum-disagreement gate can reject comparisons with little
room for a meaningful difference before answers are opened. Passing the
gate gives no evidence of superiority, calibration, or significance.
Equations~\eqref{eq:decomposition}--\eqref{eq:disagreement} are elementary
bookkeeping tools; we claim no new theorem or population guarantee.

\paragraph{The estimand depends on the control.}
For two policies on the \emph{same} recorded candidates, define the paired
query contrast $\delta_q(g,t)=H_q(g)-H_q(t)\in\{-1,0,1\}$. Its observed
average is $(F-B)/N$, rather than the difference between two independently
sampled success rates. In the matched H and I ablations, $t$ is a
separately fitted ranker with graph columns withheld or shuffled. Thus a
positive contrast establishes an improvement relative to that specified
\emph{training and feature intervention} on the observed queries. It does
not identify a causal effect of access to a graph data structure: the
relation facts are absent from the text-only control, while a same-fact
text representation is a different intervention. Nor does the paired
contrast identify the marginal value of adding graph access to a deployed
retrieval pipeline, because the shared candidate set and upstream calls
are held fixed by construction.

\paragraph{Why stronger conclusions do not follow.}
Consider any recorded fact set $F_q$ and a lossless serialization $s(F_q)$.
The composition $g\circ s^{-1}$ returns exactly the choices of a graph
policy $g$ wherever the serialization is parseable; this supplies a
counterexample to a claim that the same facts \emph{must} be represented
as a graph to produce those choices. It does not supply an efficient
learner: a bounded model might fail to parse $s(F_q)$, and generating the
serialization may itself require graph traversal. Conversely, a graph
feature model can gain over a fact-deprived text ranker yet lose to a
support-first graph rule, because equal inputs impose no ordering on
different policies' decisions. The two frozen cohorts exhibit precisely
that logical configuration. Resource parity, generalization and human
product relevance each require their own measurement, not an inference
from the paired hit counts.

\section{Frozen results}
\paragraph{What was frozen.}
The H method, 384-query sample, query-level denominator, ten-comparison
family and response parsing policy were fixed before H query preparation.
The development rule required positive gains over ten controls, at least
five distinct fixes over RRF and no development sub-cohort RRF loss worse
than two. That gate selected whether to run H, not a test statistic.
The I method was separately frozen after S fitting and development, with
768 evaluation queries and a \emph{different} 15-comparison family. Neither
evaluation is an untouched official TEST set; previous experiments may
have influenced the selected methods. Within each cohort all compared
policies reuse the available rankings and fixed candidate pools. A strict
ranking must contain every candidate exactly once. When one service fails
parsing, the valid ranking is reused for all fusion methods; when both
fail, the query remains in the denominator with zero fusion hits.
\begin{table}[t]
\centering
\caption{Hit@1 counts on disjoint official TRAIN cohorts. H evaluates the
18-feature fusion; I evaluates a distinct adapted 20-feature fusion. All
methods share each cohort's fixed 20-candidate pool. These are not repeated
tests of a single method.}
\label{tab:main}
\small
\begin{tabular}{lrr}
\toprule
Policy & H / 384 & I / 768 \\
\midrule
Two-LLM RRF & 261 & 550 \\
Max-world support & 276 & 570 \\
Top-link support & \textbf{279} & \textbf{573} \\
Text fusion without relation facts & 264 & 549 \\
Shuffled-graph fusion & 263 & 552 \\
Graph fusion (18D on H; adapted 20D on I) & 278 & 568 \\
\bottomrule
\end{tabular}
\end{table}

On H, the graph method exceeds matched text by 14 hits and shuffled graph by
15. One-sided exact paired discordance tests were adjusted by Holm over the
\emph{full ten} prespecified comparisons; adjusted $p$ values are 0.002594
and 0.003279, respectively. The corresponding adjustment for RRF (+17 hits)
is 0.023952; Borda makes the same choices as RRF here, not an independent
replication. The graph ranker has 13 fixes and 14 harms against top-link (net $-1$,
adjusted $p=1$ for graph superiority). On I, the adapted graph exceeds its
matched text by 19 hits (24 fixes, 5 harms; adjusted $p=0.003277$) and the
shuffled control by 16 (21 fixes, 5 harms; adjusted $p=0.013717$), with Holm
applied to \emph{all 15} prespecified comparisons. Against top-link the adapted
method instead has 20 fixes and 25 harms (net $-5$; adjusted $p=1$). The
adapted-vs-original graph difference on I is +5, not significant after the
prespecified adjustment. These families do not correct for all method
selection in the preceding research history; nonsignificance against the
rule establishes neither equivalence nor noninferiority.

\begin{table}[t]
\centering
\caption{Selected paired results for the frozen primary method versus each
control. The Holm families contain all 10 H and 15 I prespecified comparisons,
not only the displayed rows.}
\label{tab:paired}
\small
\begin{tabular}{llrrr}
\toprule
Cohort & Comparator & Fix / harm & Net & Holm $p$ \\
\midrule
H & Matched fact-deprived text & 15 / 1 & +14 & 0.002594 \\
H & Candidate-shuffled graph & 17 / 2 & +15 & 0.003279 \\
H & Two-LLM RRF & 26 / 9 & +17 & 0.023952 \\
H & Top-link support & 13 / 14 & $-1$ & 1 \\
I & Matched adapted text & 24 / 5 & +19 & 0.003277 \\
I & Adapted shuffled graph & 21 / 5 & +16 & 0.013717 \\
I & Two-LLM RRF & 37 / 19 & +18 & 0.111207 \\
I & Top-link support & 20 / 25 & $-5$ & 1 \\
\bottomrule
\end{tabular}
\end{table}

On H, 20,000 paired-query bootstrap resamples gave Bonferroni-adjusted
approximate percentage-point intervals of $[1.042,6.771]$ relative to matched
text and $[-4.167,3.646]$ relative to top-link. On I the analogous intervals
are $[0.5208,4.6875]$ and $[-3.2552,1.8229]$. Resampling assumes an
appropriate query-sampling interpretation and is not a finite-population
guarantee. Holm controls each fixed within-cohort family, not every adaptive
choice made earlier in the research history.

For a fixed comparison of $g$ against $t$, $F$ is the number of
$g$-only hits and $B$ the number of $t$-only hits. Under an exchangeable
null on the $F+B$ discordant queries, its one-sided exact tail is
\begin{equation}\label{eq:exact}
 p_{g>t}=\sum_{k=F}^{F+B}\binom{F+B}{k}2^{-(F+B)},
 \quad p_{g>t}=1\ \text{if }F+B=0.
\end{equation}
Holm adjustment was applied to every raw test in each prespecified family,
including comparisons not printed in Table~\ref{tab:paired}. Paired bootstrap
resamples share the same query indices across all arms (20,000 draws per
cohort); approximate Bonferroni simultaneous 95\% intervals are reported
where relevant. Neither computation removes dependence induced by the
earlier iterative research process.

The frozen 20-item pools contain labeled answers for 354/384 H queries and
704/768 I queries, but this label oracle is not a deployed policy. H required
384 shared plans, 768 ranking requests and 7,680 local Qwen evaluations;
three rankings failed strict parsing and remain in the denominator. I had
three strict ranking parse failures among 1,536 ranking requests, also
retained. These are query-level evaluations, not independent response-level
samples or a matched end-to-end cost comparison.

\section{Equal-fact alternatives and decision opportunities}
After H and I became exposed, we serialized the cached reference options,
their weights, typed and directed edges, both rankings, and candidate order
as parseable text. A separate label-blind parser reproduced \emph{all} top-link
choices on H (384/384) and I (768/768). Independent scalar reconstruction of
the six original relation columns matched 46,080 H and 92,160 I stored
numbers within $10^{-10}$. This establishes that these particular facts and
the rule can be implemented from textual serialization; relabeling the same
columns as ``text'' would not establish an alternative learner.

We therefore also tested an H-informed, post hoc structured-text learner:
18 lossy token statistics parsed from the same cached relation facts, combined
with the original 12 semantic features, fitted with the original 768 query
labels and LambdaRank setup. On \emph{already exposed H}, it hits 278/384,
the same as graph fusion (4 fixes, 4 harms; 358 identical choices), while
top-link hits 279/384. This is a development result, not proof of equivalence,
equal costs, an equal-token LLM text baseline, or independent validation.
For exploratory MAG author retrieval, the author-field text join copied the
graph path's decisions on linkable TRAIN queries. MAG has neither an independent
task evaluation here nor a transferable superiority result.

Identical action policies cannot differ in Hit@1 under any answer labels.
Before scoring a new comparison, label-blind action disagreement is therefore
a necessary check for decision opportunity, not a significance test. A frozen
384-ID prospective Amazon TRAIN queue exists, but no queries in that queue
have generated scored actions: the launch preflight stopped before query-text
access and external calls. We do not include it in any result table.

\paragraph{Where are the gains?}
An exploratory analysis defined graph activation after observing H/I as a
maximum-support domain smaller than 20. Activation occurs for 55/384 H and
94/768 I queries. On H, the 18-feature graph fusion fixes 13 and harms zero
relative to text within the 55 activated queries; the remaining queries
account for a net gain of one. On I, the adapted method fixes 17 and harms
one relative to its adapted text arm within the 94 activated queries, with
a net gain of three outside. This accounting describes the already exposed
samples; it is not a prespecified subgroup test or proof that a particular
edge caused a successful prediction. On the same activated sets, top-link
hits 37/55 H and 67/94 I, versus the respective fusion methods' 32/55 and
59/94. An informative activated subset can coexist with an inferior learned
decision policy.

\paragraph{Why is the rule hard to beat?}
Only 145 of the 768 queries used to fit the original fusion have a reference
mention; 111 have nonconstant graph features, and 58 have both labeled
positive and negative candidates within the maximum-support domain. On
exposed I, the adapted fusion and top-link choose different products on
197 queries, but only 45 choices change Hit@1 status (20 favorable to fusion,
25 favorable to the rule). These counts are consistent with scarce
decision-relevant supervision and harmful overrides, but do not isolate
training coverage from representation, support-first policy, or distribution
effects. They are not an additional evaluation set.

\section{Limitations and research status}
\paragraph{How the three questions interact.}
An ablation comparison holds supervision, candidate pool and much of the
scoring stack constant while removing graph-derived columns. Its positive
paired outcome identifies a benefit of that feature block \emph{against
that fact-deprived arm} on the frozen queries. It does not identify
whether the information can be encoded as text, whether an LLM can use
the encoded facts under a fixed token budget, or whether a different
decision policy would have done better with the same facts. The rule
replay shows constructively that textual serialization can reproduce a
specific deterministic decision. The H-informed text learner goes beyond
relabeling graph columns, but is post hoc and still not a prospective
equal-resource study. Finally, top-link is a graph-using rule; its higher
hit count is a failure of the \emph{learned policy's superiority claim},
not evidence that graph relations are useless. This separation is the
empirical point of the paper: the strongest observation and the strongest
counterexample address different questions.

\paragraph{Mechanistic interpretation is conditional.}
Under the observed link options and fixed pool, a support domain often
retains a labeled answer. But a graph edge is merely an existence proxy:
related accessories, alternate products and incompatible items can all
be linked. The I action accounting makes the bottleneck concrete. The
adapted learner and top-link disagree on 197/768 selected candidates,
yet their benchmark hit states differ on only 45 queries. To beat the
rule by selecting the learner selectively, one would need to recover
enough of its 20 unique fixes without triggering its 25 unique harms;
this is a sample accounting condition, not an identified way to build
such a selector. Enlarging candidates, improving reference linking,
training a role-aware score and imposing a support-first constraint
would modify different components. None has been shown in this study
to reverse the frozen strong-rule comparison under new evaluation.

These comparisons isolate an information increment against \emph{fact-deprived}
text, not against a trained text system with identical relation facts and
matched indexing, latency, peak memory and price. Post hoc cached-feature
and paired construction microbenchmarks exclude retrieval, model loading,
reference API requests and Qwen; they cannot establish deployment efficiency.
Top-link uses graph facts too, and its stronger outcomes do not mean that a
text-only retriever dominates graph retrieval. The observations are about
fixed-pool selection, not full-corpus end-to-end STaRK performance.

H/I are exposed TRAIN cohorts under distinct prespecified families, not
independent final confirmation; synthetic qrels are not human semantic
judgments. A prepared blind review of H/I decisions has no completed human
labels. The new cohort remains unscored pending enforceable service spending
limits, frozen matched-resource protocol, source and provenance checks, and
the label-blind action gate. Future work must accept negative outcomes,
assess robustness beyond a single shuffle seed, compare equal-fact methods
on new queries with matched end-to-end resources, and complete human review.
The current evidence archive is locally auditable but is not yet an approved,
redistributable public reproduction package.

\section{Conclusion}
Relational features add value over text without relation facts in the two
studied Amazon settings, but that increment does not establish a uniquely
necessary graph representation or a learned decision advantage over the
shared-input top-link rule. The three claims require distinct controls.
Until new-cohort scoring, matched end-to-end comparison and semantic review
are complete, this is a bounded empirical working draft rather than an
independently confirmed or submission-ready superiority claim.

\bibliographystyle{plainnat}
\bibliography{references}

\appendix
\section{Evidence roles and reproducibility notes}
The present manuscript is a condensed English review draft of the current
Chinese working paper, not a second experimental result. Table~\ref{tab:roles}
prevents query counts from being reused as independent confirmations. Earlier
validation experiments tested different repair policies and are not pooled
with H/I. All numerical claims in Tables~\ref{tab:main} and~\ref{tab:paired}
refer to archived frozen actions and query-level results. Independent readers
need access to those records and authorized data before claiming a full
reproduction; a local consistency check is not independent retraining.

\begin{table}[h]
\centering
\caption{Data-use roles. Counts refer to queries, not LLM calls or candidate
products. No official TEST or human-query answers are evaluated.}
\label{tab:roles}
\small
\begin{tabular}{lrl}
\toprule
Set & Queries & Use and evidential status \\
\midrule
A/B/F (TRAIN) & 768 & Original fusion fitting \\
C/D/E (TRAIN) & 384 & Repeated development; exposed \\
H (TRAIN) & 384 & Frozen 18D evaluation; now exposed \\
S (TRAIN) & 1,024 & Adapted semantic encoder training \\
I (TRAIN) & 768 & Distinct frozen 20D evaluation; now exposed \\
Prospective ID queue (TRAIN) & 384 & IDs frozen; no scored actions \\
Official TEST / human & 1,642 / 81 & Not used for reported results \\
\bottomrule
\end{tabular}
\end{table}

The model comparison was fixed before H query preparation, with 384 plans
and 768 rankings budgeted; the adapted comparison on I specified 768 plans
and 1,536 rankings. Strict parsing failures stay in denominators, with
available ranks reused consistently across same-cohort arms. The exact
one-sided paired tests compare fixes to harms across query IDs, with Holm
adjustment within each frozen family. A same-result action is never an
additional successful replication. The structured-text arm, activation
subgroups, and the MAG author-field join are explicitly post hoc or TRAIN
developmental; none enter the original H/I comparison families. The staged
prospective ID queue does not permit retroactive selection or scoring. Its
zero-call preflight and lack of a metered cost stop prevent us from claiming
that confirmation has started.

Publication requires separate review of style-year rules, author approvals,
data/service-output redistribution rights, artifact privacy, checklist
answers and the eventual scientific claim. Anonymous styling and local
typesetting do not constitute a conference submission.
\section{Evidence and completion status}
This appendix separates protocol details, equal-fact comparisons, local
resource accounting, exploratory MAG observations and the complete frozen
Amazon comparison families. Each has a different evidential role: a
reconstructed action is not a new labeled evaluation, a local timing is
not a deployed cost, and a frozen comparison family does not account
for earlier choices in the research program.

\section{Claim-to-evidence trace and reproducibility boundary}\label{app:trace}
The local analysis tree contains archived protocols, request and action
snapshots, source hashes and validation reports. A reader can verify
arithmetic without rerunning models, but cannot independently reproduce
end-to-end behavior from the PDF alone. Table~\ref{tab:trace} specifies
what evidence should be inspected for each claim. Paths are workspace
relative; access to a path is not a license to redistribute its contents.
This manuscript's English wording must remain consistent with the
more detailed Chinese working draft and the read-only evidence validator.

\begin{table}[h]
\centering
\caption{Audit targets under \texttt{output/analysis/}. Numeric IDs identify
the corresponding numbered report; H, I and queue IDs denote frozen
experiment directories. A local validation is not external replication.}
\label{tab:trace}\small
\begin{tabular}{p{5cm}p{2.7cm}p{4.5cm}}
\toprule
Claim & Evidence ID & What it checks \\
\midrule
H 18D comparison & H 20260922 & Protocol, actions, 384 hits, ten tests \\
I 20D comparison & I 20260923 & Different method, 768 hits, 15 tests \\
Same-fact rule replay & Report 66 & H/I label-blind action identity \\
Six graph columns & Report 67 & Scalar tolerance and input hashes \\
H text learner & Report 68 & Post hoc fit, gate, actions \\
Local resource scopes & Report 70 & Separate timed boundaries \\
Unscored ID queue & Queue v0 & Zero-access original protocol \\
Stopped launch & Launch v3 & No query/call clearance \\
Human review design & Report 65 & Prepared rubric, zero labels \\
\bottomrule
\end{tabular}
\end{table}

H and I are the directories \texttt{multiview\_transfer384\_20260922/}
and \texttt{disjoint\_transfer768\_20260923/}. Numbered IDs refer to
the matching Markdown report prefix in the analysis directory. Queue v0
and Launch v3 refer to \texttt{amazon\_candidate\_cohort\_20260925\_v0/}
and \texttt{amazon\_prospective\_launch\_20260925\_v3/}, respectively.

The read-only command \texttt{perl -w paper\_evidence\_check.pl}, from the
repository root, checks H/I summary and paired rows against frozen result
files, original source and input hashes for the H text pilot, candidate-ID
disjointness and zero-access status, and the two distinct H resource
benchmarks. It does not fetch the official dataset, rebuild the BM25/E5
indices, reproduce service outputs, retrain Qwen, run a blinded reviewer,
or prove that outside access to the candidate suffix never occurred.
Some raw records and model files are locally archived under an ignored
analysis tree; the project has not delivered a cleared public package.
The generated PDF and all third-party styles are reading artifacts, not
a substitute for inspectable authorized sources and exact version pins.

\paragraph{What a third-party reproduction would have to specify.}
At minimum: the STaRK revision and splits, full-text and graph schema,
hashes of every source/input/model, the lexical query tokenizer,
BM25/E5 initial retrieval and tie ordering, exact prompt and model IDs,
per-request caps and error semantics, entity-linking options/weights,
Qwen scoring revision, fit IDs and hyperparameters, 20 candidate IDs for
every query, action serialization, and a rule for missing responses.
A new run must specify whether API usage is billed and which units are
returned by the provider. A local syntax or checksum check does not
validate the truth of the benchmark labels or the legitimacy of
redistributing the underlying responses. Independent execution may
produce different API responses if the pinned service implementation
is unavailable; in that case the difference must be reported rather
than silently substituted with another model.

\section{Human semantic review is a separate estimand}\label{app:human}
The official synthetic qrels measure agreement with the benchmark
answer construction, not independent assessment of product role,
compatibility or intended purchase. The already prepared H/I blind
review draws per cohort 48 graph-versus-top-link \emph{disagreement}
queries and 24 queries on which they agree: 144 tasks in all. H has
73/384 differing actions and I has 197/768; the sample is chosen from
these exposed TRAIN cohorts, so any eventual semantic result remains
post hoc and cannot transform them into an untouched confirmation.
Each task displays at most the two products selected by the studied
policies. It does not expose all 20 retrieved products and cannot
estimate full-pool human recall.

Two real independent reviewers should receive separately shuffled
blinded forms containing the same source query, chosen-product text,
tags and recorded reference-edge facts. They should not be shown
policy names, frozen benchmark labels, previously computed outcomes
or the private mapping from displayed options to policy identities.
Each must label product role match, explicit textual conflict,
relational-condition support and overall suitability for \emph{each}
displayed option, using yes/no/uncertain with a cited justification.
An also-buy/also-view edge is evidence of co-occurrence, not proof
of semantic compatibility; missing graph evidence is not a verified
negative. Separating the two reviewers' access from the workspace
containing the private mapping is an operational requirement, not
something that a validator script can ensure.

Before unblinding, validate completeness and the two identities, archive
first-pass judgments, report uncertain rates and categorical agreement,
and perform a separately recorded adjudication. Only then join the
graph-versus-rule actions to suitability judgments, reporting fixes,
harms, both-yes, both-no and unresolved separately for H and I. Treat
unresolved paired differences conservatively with their possible
$[-1,1]$ contribution rather than dropping them. The agree-action
tasks serve as rubric and baseline checks, not evidence of a policy
difference. The disagreement-stratified sample is not IID from
all benchmark queries. These are analysis instructions; as of this
draft \emph{no human annotations have been completed}. Even completed
judgments cannot replace a new, independently frozen algorithmic
evaluation.

\section{Prospective study: stop conditions, not outcomes}\label{app:future}
A later preflight froze 384 official TRAIN \emph{IDs} selected from the
sorted-ID permutation at positions 4672--5055. Its immutable v0 record
sets paid calls to zero and forbids opening the new query text and
answers. It proposes a graph policy, the separately H-designed
structured-text learner and top-link on identical 20-candidate pools,
with a label-blind minimum net action-disagreement threshold of five
against \emph{both} controls. This gate is necessary for a meaningful
Hit@1 contrast but cannot predict the sign of the contrast. Passing
it would not retrospectively remove H's influence on text-arm design.

User authorization later specified a task-bounded maximum of 1,152
external requests, conditional on a complete protocol and on the
action gate before reading answers. The archived executor caps the
attempt count and per-request output tokens and verifies observed
service IDs, but records input usage and a provider cost field only
\emph{after} each call. It does not enforce a cumulative monetary
hard stop before dispatch, and the cost field's billing unit and an
explicit finite monetary ceiling remain unresolved. The separate
launch preflight therefore returned STOP: no new-query text, answers
or paid models were accessed by that attempt. Neither a user's
allowance estimate nor a locally frozen ID list is proof that the
cost-control or provenance conditions have been met.

To turn that queue into a real prospective experiment, one must
first resolve remaining possible exposure outside known local
metadata, verify the original service models without disclosing
credentials, bind the action producer and its input/model files to
an immutable protocol, implement enforceable token and monetary
limits, and specify abstention and missing-call treatment. Only
then may the approved query-only action generation run. The
independently fitted equal-fact text arm should use the same cached
edge/link facts, candidate pool and supervision and declare whether
the target is pure quality at fixed facts or quality per measured
resource. If full cold/warm lookup, comparable index construction
and updates, token counts, latency, RAM/VRAM and amortized costs
cannot be measured, that run cannot establish deployment superiority.
Scoring any answers remains conditional on an archived label-blind
action file that passes both predeclared gates. An unfavorable outcome
must be reported without extending the queue or changing the policy.

This leaves four independent completion criteria for a mature paper:
(i) an eligible prospectively scored comparison under correct
model/budget guards, (ii) genuinely matched end-to-end resource
measurements or an explicitly quality-only claim, (iii) actual
independent human semantic judgments with reviewer separation, and
(iv) author-approved reproduction rights and a redistributable
evidence bundle. None of these is manufactured by typesetting.

 \section{Implementation, data and frozen decisions}\label{app:implementation}
 \paragraph{Dataset and candidate universe.}
 The local STaRK-Amazon snapshot contains 957,192 products, 1,035,542
 nodes, and 5,382,103 stored directed edges. Stored edge counts depend on
 the snapshot's counting convention; they need not match a graph statistic
 reported for another revision. The full graph contains also-buy,
 also-view, brand, category and color edges, but the support proxy in this
 paper uses the first two types. The official synthetic query split contains
 5,910 TRAIN, 1,548 validation and 1,642 TEST queries, plus 81 human queries.
 Synthetic query generation uses both text and graph information; that
 design may favor relational features relative to naturally sampled tasks.
 Product texts, graph structure, queries and model pretraining corpora can
 overlap even when evaluation \emph{query IDs} are disjoint. A product not
 in the labeled answer set is a benchmark non-hit, not necessarily a human
 verified semantic negative.

 \paragraph{Retrieval, ranking and failure handling.}
 Frozen full-corpus BM25 and E5 ranks are joined using RRF with rank
 constant 60, and the first 20 candidates are retained. This candidate
 retrieval rank is distinct from the two subsequent LLM ranks of those
 20 candidates. Both the learned fusion and the top-link rule use the same
 frozen candidate order as the ultimate tie-breaker. The two service IDs
 recorded for H are \texttt{gpt-5.6-sol} and \texttt{gpt-6-astra}; the
 identity check uses reported service usage rather than model self-report.
 The reference plan is a separate shared request to the latter service,
 limited to three literal product mentions in the query. An exact title
 match is attempted before dense E5 link proposals. Each LLM sees at most
 512 \texttt{o200k\_base} tokens of product text plus capped tags, but
 the E5 retrieval truncation and Qwen effective contexts are not all
 identical token windows. No claim of equal input length across encoders
 follows from sharing product identities.

 A valid ranking contains every candidate exactly once. No ranking is
 silently repaired by appending absent products, and no parsed failure is
 removed from the query denominator. All arms in a cohort use the same
 usable service outputs. H had valid reference plans for all 384 queries
 and 765 valid ranking responses out of 768; I had valid plans for all 768
 and 1,533 valid rankings out of 1,536. The original 18D fusion on H
 scores 7,680 product/query pairs locally with a frozen Qwen3-Reranker-0.6B
 yes-versus-no logit difference. Calls to obtain ranks and reference plans
 are shared by top-link, but the rule does not require Qwen prediction or
 learned-model inference to make its decision. A same-query prediction
 cost measured after all upstream features are cached cannot be used as
 its deployment cost.

 \paragraph{Six rank and six nonrelational columns.}
 For each valid model $m$, write $r_m(v)$ for its rank and $M=|\mathcal M_q|$.
 The rank block actually used on the fixed $K=20$ pool is
 \begin{align}
 \phi_R(q,v)=\big[& (\min_m r_m(v)-1)/19,
  (\max_m r_m(v)-1)/19,\\
 & (M^{-1}\!\sum_m r_m(v)-1)/19,
  (\max_m r_m(v)-\min_m r_m(v))/19,\nonumber\\
 & M^{-1}\!\sum_m 1/r_m(v),
  M^{-1}\!\sum_m\mathbf{1}\{r_m(v)=1\}\big].\nonumber
 \end{align}
 The nonrelational block holds normalized initial retrieval rank;
 query/title and query/text term overlaps; a capped normalized product-text
 length; the frozen Qwen relevance logit difference $z(q,v)$; and
 $z(q,v)-\max_{u\in C_q}z(q,u)$. The Qwen instruction distinguishes a
 requested purchase from a reference item already owned by a user, but
 its score is not a verified role or calibrated probability. The text
 ablation is $[\phi_R;\phi_T]$, the graph method appends $\phi_G$,
 and shuffle preserves the six graph columns' marginal values while
 moving the candidate association as one block within each query.
 Shuffling is controlled by a fixed seed and query identifier. Only one
 prespecified shuffle realization was evaluated; even a successful
 paired comparison does not estimate variance across many shuffles.

 \paragraph{Six graph columns and policy.}
 For options $R_{qj}$, the similarity softmax is
 \begin{equation}
  w_{qjt}=\frac{\exp(s_{qjt}/0.05)}
  {\sum_{u\in R_{qj}}\exp(s_{qju}/0.05)},\qquad
  p_q(h)=\prod_jw_{qj,h_j}.
 \end{equation}
 Exact-title option ties instead receive uniform weights. If $A(v,t)$
 denotes the observed existence of an also-buy or also-view edge in
 either direction, $U_q(v,h)=J_q^{-1}\sum_j A(v,h_j)$ for nonzero $J_q$.
 We then take $\phi_G=[\mathbb E_{\rm unif} U,\min U,\max U,
 \operatorname{Std}_{\rm unif}U,\mathbb E_{p_q}U,
 \operatorname{Std}_{p_q}U]$. These moments flatten typed directed
 records into a coarse undirected existence indicator for scoring.
 The fact serialization used by the equal-fact controls retains the
 underlying edge types, directions and duplicates instead. This distinction
 matters: the six coarse values are not a semantics oracle. If no mention
 is identified, the whole graph block is zero.

 For top-link let $h_j^*=\arg\max_{t\in R_{qj}}w_{qjt}$, resolving a link
 tie by frozen option order. Its lexicographic action is
 \begin{equation}
 \pi_{\rm top}(q)=\arg\max_{v\in C_q}^{\rm lex}
  \left(U_q(v,h^*),\sum_{m\in\mathcal M_q}\frac{1}{60+r_m(v)},
  -\operatorname{rank}_{C_q}(v)\right).
 \end{equation}
 The second and third entries matter only after support ties. This is
 not a ``weak text-only'' baseline: it shares links and graph data with
 the graph fusion, but uses a different decision constraint. Max-world
 support replaces the first entry by $\max_hU_q(v,h)$, with the same
 two-rank tie-breaker. Its 276 H and 570 I hits are below top-link but
 above either learned graph method on I. A model trained without the
 top-link support column cannot be declared to have learned the
 support-first policy merely because the underlying facts were available.

 \subsection{Equal-fact controls: what is and is not tested}\label{app:equalfacts}
 \paragraph{Rule and feature replays.}
 The post hoc rule control serializes each query's frozen candidate order,
 rankings, reference option IDs, option weights, and all recorded connected
 candidate/option edge IDs, types, directions and duplicates as UTF-8
 structured text. An independent parser then selects each top link,
 recomputes its support and applies the same lexicographic RRF tie-break.
 It agrees on 384/384 H actions and 768/768 I actions; only 68 H and 138 I
 queries contain a named reference. The serialized per-query size spans
 215--8,480 H bytes and 228--3,891 I bytes. These counts specify \emph{bytes},
 not a token budget for a model. A second label-blind replay enumerates
 all linked-option worlds and reconstructs the original graph block
 to absolute error at most $10^{-10}$ for every one of the 46,080 H and
 92,160 I scalars (queries times 20 candidates times six columns).
 That construction can be performed in either a graph-oriented or
 text-parsed program. It does not establish that an off-the-shelf
 language model, given the same tokens, will learn an equivalent policy.

 \paragraph{An H-informed learned text arm.}
 To avoid calling a mere reserialization of six graph columns a learned
 text baseline, a separate H development experiment feeds the losslessly
 serialized \emph{facts} to a parser that constructs 18 lossy lexical
 statistics per candidate. These are combined with the original 12
 nonrelational features and fitted with the same 768 fitting queries,
 query groups, labels, LambdaRank settings and tree count. Original text
 and original graph models were retrained in the same CPU session and
 replayed their frozen H actions. The new text model makes 278/384
 H hits, against 278/384 graph and 279/384 top-link. Relative to graph,
 its actions agree on 358 queries; of 26 disagreements only eight change
 benchmark hit status, four in each direction. Relative to the rule
 there are 13 fixes and 14 harms. The fit IDs and H query IDs are disjoint,
 but the design of this text arm was chosen after H outcomes had already
 been inspected. A tie in this exposed-cohort outcome is not evidence
 of equivalence or a matched-cost performance guarantee. The parser is
 lossy despite lossless input serialization, and it is neither an
 LLM-baseline evaluation nor an adapted-I semantic-encoder control.

 \paragraph{Why input equality has several meanings.}
 We distinguish the underlying fact set, its textual or graph-oriented
 representation, model supervision, candidate pool, upstream ranking
 calls, and measured deployment resources. The original text12 ablation
 shares a candidate pool, rankings, semantic scores and label budget,
 but does \emph{not} have the link and edge facts. The post hoc parser
 has those facts, but is a deterministic program rather than a matched
 token/latency LLM. The H structured-text learner has the facts and
 original fit labels, but was designed after exposure and uses lossy
 lexicalization; there is no corresponding I adapted-text same-fact
 learner. The top-link rule has underlying facts but differs in features
 and support-first action logic. Conflating these controls would turn
 one valid ablation result into an invalid graph-exclusivity claim.

 \subsection{Measured local resources versus deployment cost}\label{app:resources}
 Two separate exploratory benchmarks reused \emph{exposed H}; neither read
 new query answers or assessed quality on a new cohort. The first used
 seven serial repeated predictions over 384 queries and 7,680 candidates,
 alternating graph and text arms with one prediction thread. Its graph
 input features had already been cached. The text arm separately
 serialized relation facts into JSON and parsed its 18 columns before
 model inference. Both models reproduced their original H actions on
 each repeat. Table~\ref{tab:resource} separates these stages; timings
 are descriptive local measurements, not confidence intervals or
 end-to-end deployment comparisons.

 \begin{table}[h]
 \centering
 \caption{Exposed-H local resource observations, seven runs per stage.
 Different protocols must not be compared across the panel boundary.}
 \label{tab:resource}\small
 \begin{tabular}{p{8.7cm}rr}
 \toprule
 Stage and measurement boundary & Median (s) & Range (s) \\
 \midrule
 Text JSON serialization and parse; no index & 0.550701 & 0.547480--0.756351 \\
 Graph model prediction; graph features cached & 0.023100 & 0.022765--0.030991 \\
 Text model prediction; text features parsed & 0.023335 & 0.023135--0.031144 \\
 \midrule
 Graph six-column construction, \emph{preloaded} facts & 0.126542 & 0.125997--0.128655 \\
 Text JSON and 18-column construction, \emph{same preloaded} facts & 0.113578 & 0.112287--0.114158 \\
 \bottomrule
 \end{tabular}
 \end{table}

 The second benchmark independently and symmetrically constructs graph
 and structured-text columns from the \emph{same preloaded} H case and
 link-option data, seven serial repeats with alternating arm order. Its
 graph producer exactly recovers the 46,080 frozen graph scalars on every
 repeat. The first benchmark includes Python \texttt{tracemalloc} while
 timing text serialization; it also records a 4,311,096-byte median
 Python allocation peak during that parse, \emph{not} process RSS.
 It is invalid to compare its 0.550701-second text parse directly with
 the other run's 0.126542-second graph construction as an efficiency
 ratio. The construction outputs also have different feature dimensions.

 Both benchmarks omit full-corpus index build/update, BM25/E5 retrieval,
 LLM reference generation and two ranking requests, entity-linking
 execution, Qwen scoring, model loading, process/VRAM peak and controlled
 cold/warm end-to-end latency. The recorded service \texttt{cost} fields
 have unverified units; dollars cannot be reconstructed reliably from
 them. Machine resources were not exclusively reserved to certify
 same-time-budget deployment. Neither arm is shown to be cheaper or
 noninferior under an actual matched end-to-end resource budget.

 \subsection{Exploratory cross-task and negative-control evidence}\label{app:mag}
 The exploratory STaRK-MAG author task retrieves papers, not Amazon
 products, and neither Amazon model is transferred to it. On a selected
 1,024-query MAG TRAIN development sample, author-writes reranking
 hits 376, author-enhanced BM25 363 and a mismatched-support control
 299. Yet exact matching against the same index's author text field
 reproduces every graph-rule action on its 122 linkable queries. This
 supports a counterexample to a \emph{representation-exclusive}
 interpretation of that particular graph path; it is not independent
 cross-dataset confirmation of Amazon's learned fusion.

 Other MAG TRAIN pilots supply deliberately negative diagnostic evidence.
 On 259 disjoint TRAIN queries, an external-author-ID graph profile hits
 168, exact-name text 185 and BM25 153; these graph and text profiles are
 \emph{not} strict equal-fact representations. On 434 collaborator
 queries a graph path and its named-text comparator differ on only four
 top-1 actions. On 512 institution queries, graph and same-fact text
 prompts differ on only nine actions. Those latter two comparisons failed
 their prespecified label-blind disagreement gates; their answers were
 not opened and no hit rates for them are claimed. The prompts do not
 have exactly equal lengths even when they share facts and model access.
 By Eq.~\eqref{eq:disagreement}, any Hit@1 difference on their fixed
 cohorts is at most $4/434$ or $9/512$ respectively, regardless of
 which policy would benefit. We did not evaluate MAG official VAL/TEST
 or human queries for these paths.

\clearpage
\section{Full frozen comparison families}\label{app:families}
Tables~\ref{tab:fullh} and~\ref{tab:fulli} show every prespecified
comparator, not only the favorable rows in the main paper. The first column
identifies a different policy on the \emph{same} cohort: one model's raw
ranking, a single-model LambdaRank (LTR), a linked-gain single-model
postprocessor, rank aggregation, or a graph-support rule. These policies
can differ in minimal deployment cost even when evaluated on the same
cached candidates and service outputs. In particular, rank-only policies
need not run the extra Qwen scorer or reference-extraction path if deployed
independently. Each table uses a \emph{single} numerator (the frozen graph
method) and a query denominator fixed before result inspection. A negative
net count is a result, not a missing experiment.
\begin{table}[h]
\centering
\caption{Complete H family ($N=384$; primary: 18D graph fusion, 278 hits).
Fix/harm is primary-only hit/comparator-only hit. Adjusted $p$ uses all ten
one-sided paired comparisons; rounding is for display only. LTR and linked
postprocessors use the indicated single ranking.}
\label{tab:fullh}\small
\begin{tabular}{lrrrl}
\toprule
Comparator & Hits & Fix/harm & Net & Holm $p$ \\
\midrule
Text fusion (12D) & 264 & 15/1 & +14 & 0.002594 \\
Shuffled fusion (18D) & 263 & 17/2 & +15 & 0.003279 \\
RRF & 261 & 26/9 & +17 & 0.023952 \\
Borda & 261 & 26/9 & +17 & 0.023952 \\
Single-model LTR 0 & 270 & 18/10 & +8 & 0.554800 \\
Single-model LTR 1 & 273 & 15/10 & +5 & 1 \\
Single-model linked 0 & 273 & 17/12 & +5 & 1 \\
Single-model linked 1 & 277 & 14/13 & +1 & 1 \\
Max-world support & 276 & 14/12 & +2 & 1 \\
Top-link support & 279 & 13/14 & $-1$ & 1 \\
\bottomrule
\end{tabular}
\end{table}

The H training pool consists of A/B/F, not H labels; development on C/D/E
was used to choose the method. The 384 H queries were held out \emph{at
the time of that freeze}, not across the complete research history. RRF
and Borda coincide in this cohort, making their two rows two controls in
the fixed family but not independent empirical replications. The apparent
positive difference over maximum-world support is only two queries. The
top-link control beats the main graph method by one hit; its nonsignificant
comparison must not be reversed into an equivalence claim. Complete
query-level action and hit records reside under the frozen H results
archive (see Appendix~\ref{app:trace}).

\begin{table}[h]
\centering
\caption{Complete I family ($N=768$; primary: adapted 20D graph fusion,
568 hits). The older 18D text/shuffle rows do \emph{not} share the adapted
encoder's additional supervision, unlike adapted text/shuffle.}
\label{tab:fulli}\small
\begin{tabular}{lrrrl}
\toprule
Comparator & Hits & Fix/harm & Net & Holm $p$ \\
\midrule
Old 18D graph fusion & 563 & 17/12 & +5 & 1 \\
Old text fusion & 542 & 33/7 & +26 & 0.000296 \\
Old shuffled fusion & 543 & 33/8 & +25 & 0.000729 \\
RRF & 550 & 37/19 & +18 & 0.111207 \\
Borda & 550 & 37/19 & +18 & 0.111207 \\
Single-model LTR 0 & 553 & 43/28 & +15 & 0.383695 \\
Single-model LTR 1 & 562 & 34/28 & +6 & 1 \\
Single-model linked 0 & 555 & 42/29 & +13 & 0.538695 \\
Single-model linked 1 & 565 & 33/30 & +3 & 1 \\
Max-world support & 570 & 24/26 & $-2$ & 1 \\
Top-link support & 573 & 20/25 & $-5$ & 1 \\
Matched adapted text (14D) & 549 & 24/5 & +19 & 0.003277 \\
Adapted shuffled graph (20D) & 552 & 21/5 & +16 & 0.013717 \\
Original adapted-encoder score & 561 & 17/10 & +7 & 0.743366 \\
Raw adapted model ranking & 431 & 167/30 & +137 & $<10^{-20}$ \\
\bottomrule
\end{tabular}
\end{table}

Five I comparisons are positive after the specified adjustment: matched
adapted text, adapted shuffle, old text, old shuffle, and the standalone
semantic score. The last three are not substitutes for an equal-supervision
graph comparison. The raw adapted model ranking is an included, much
weaker control, not evidence of strong-baseline superiority. Neither RRF
nor either graph-support rule is among the five positive adjusted
comparisons. The adapted model is not significantly better than the
old 18D graph model on I, so the additional encoder-training step itself
has no isolated confirmed gain. This family cannot be pooled with H's
ten tests or interpreted as 1,152 repeated trials of one policy.

\section{Derivations and finite-sample scope}\label{app:derivations}
\paragraph{Error accounting.}
For a policy required to select from $E_q\subseteq C_q$, the three terms in
Eq.~\eqref{eq:decomposition} are nonnegative because
$H_q(\pi)\le O_q(E_q)\le O_q(C_q)\le1$. Summing the three brackets
telescopes to $1-H_q(\pi)$. Neither $O_q(E_q)$ nor $O_q(C_q)$ can be
computed without the benchmark labels, and neither is an attainable
deployed policy. The full-pool oracle upper bound on hit counts is 354/384
for H and 704/768 for I. Thus the 18D graph fusion's 106 H misses include
30 without any labeled answer in the fixed pool and 76 with an answer
present but not selected. The adapted fusion's 200 I misses likewise
split into 64 unavailable and 136 within-pool selection failures. This
does not identify which failures arise specifically from the graph.

\paragraph{Paired outcomes.}
For each query define $d_q=H_q(g)-H_q(t)\in\{-1,0,1\}$. Its nonzero terms
are the $F$ fixes and $B$ harms, hence
$\sum_qd_q=F-B$. If two policies return the same candidate, both hit
indicators coincide for every possible $Y_q$, so $d_q=0$ whenever
$q\notin D(g,t)$. Thus
$|\sum_qd_q|\le\sum_{q\in D}|d_q|\le |D|$, proving
Eq.~\eqref{eq:disagreement}. This bound is about the \emph{observed
fixed queries} and policies, not a rate guarantee for future queries.
The paired exact test in Eq.~\eqref{eq:exact} conditions on the observed
number of discordant outcomes. It needs the interpretation that under the
null their two signs are exchangeable; Holm's family control applies to
those particular frozen tests, not to previous model search. A confidence
interval crossing zero does not certify noninferiority. To test that claim
one would need a justified margin and a separate prespecified procedure.

\paragraph{Matched-size, matched-rank-set expectation.}
Partition each fixed candidate pool into five consecutive bins of four
by the two-LLM RRF order. In bin $b$, let $n_b=4$ candidates include
$p_b$ benchmark positives, and let a support domain select $k_b$ items.
A control set uniformly selects $k_b$ without replacement in each bin.
Exactly $\binom{n_b-p_b}{k_b}$ of the $\binom{n_b}{k_b}$ possible choices
avoid positives. Since bin draws are independent by construction,
\begin{equation}\label{eq:hypergeom}
 \mathbb E[O_q(E_{\mathrm{control}})]
 =1-\prod_{b=1}^{5}
   \frac{\binom{n_b-p_b}{k_b}}{\binom{n_b}{k_b}}.
\end{equation}
The usual convention assigns zero to a bin's avoidance probability when
$k_b>n_b-p_b$. Equation~\eqref{eq:hypergeom} is an exact expectation for
this artificial matched-set distribution, not a model of how graph edges
are generated. A more finely matched degree or within-bin rank control
might change the descriptive comparison. On the \emph{post hoc} activated
subsets ($|E_q|<20$), recorded graph-domain oracles total 43/55 H and
79/94 I; the corresponding matched-set oracle expectations sum to 25.50
and 44.52. Fractional values are expectations, not fractional queries or
additional independently sampled observations. On those same subsets
top-link hits 37 and 67, while the matched-set first-position expectations
are 16.92 and 29.67. All these numbers were computed after H/I exposure;
no new confirmatory $p$ value follows.

\paragraph{A restricted improvement ceiling.}
Consider a hypothetical new policy that differs from top-link only on
the activated set $\mathcal A$ and still selects from $E_q$ there. For
each activated $q$, $H_q(\pi)\le O_q(E_q)$, while outside $\mathcal A$
the two policies have equal hits. Consequently its \emph{observed}
improvement is bounded by
\begin{equation}
 \sum_q[H_q(\pi)-H_q(\pi_{\mathrm{top}})]
 \le\sum_{q\in\mathcal A}
 [O_q(E_q)-H_q(\pi_{\mathrm{top}})].
\end{equation}
The right-hand side is six H queries and twelve I queries, each 1.5625
percentage points of its respective full cohort. These are label-oracle
ceilings conditional on the specified domain and activation restriction;
they do not bound policies that change choices outside $\mathcal A$ or
select outside $E_q$, and they do not promise an attainable gain. For
example, the studied full-pool fusion is \emph{not} constrained by the
restriction. No novel generalization theorem is asserted here.
\end{document}